\section{软件工程应用}
\label{sec:app}

本节说明 WhatyTerm 在真实软件工程场景中的应用方式，并给出两个端到端的使用轨迹。

\subsection{典型应用场景}

\textbf{无人值守的自动化操作。} 最基础的应用是让编码代理在无人看管下持续推进任务：自动化监督层感知到代理空闲或弹出确认框时，自动注入“继续”或选择“允许本次会话”，使一次可能需要数小时、涉及数十轮交互的开发任务无需人工全程值守。

\textbf{全流程的构建—测试—交付。} 结合领域策略框架，系统可在“开发”阶段自动放行常规确认，而在“测试归档”等关键阶段强制停手等待人工确认；配合持续集成与发布流程，代理产出的改动经隔离分支提交后进入常规 CI，形成从编码到跨平台交付的链路。

\textbf{自主项目交付。} 借助三层自主循环，用户只需给出一个较高层次的需求，系统即可自动完成需求澄清、任务拆解、并行开发、集成验证与失败回退，产出一个可验证的交付物。

\subsection{案例一：交互式无人值守开发}

考虑一个“为现有 API 端点增加分页能力”的任务。用户在某会话中启动编码代理并开启自动操作后离开。此后：代理读取相关文件、提出修改方案，在需要执行文件写入时弹出权限确认——监督层依据领域策略自动选择“允许本次会话”；代理编译并运行测试，若出现瞬时 API 限流错误，系统的供应商故障转移机制自动切换到下一供应商重试；代理完成后停在空闲提示符，监督层注入“继续”推进到下一子步骤；当代理输出“任务已完成”类措辞时，死循环防护逻辑识别到无新任务，停止注入并提示用户。整个过程中，人工干预点仅为任务开始时的启动与结束时的验收。

\subsection{案例二：自主交付一个完整需求}

考虑一个更高层次的需求。系统首先在 macro 层判断需求是否需要澄清，随后生成 PRD 与带接口契约的 WBS，并将子任务按依赖拓扑排序为若干可并行批次。每个子任务进入 micro 层，由 Developer 实现、Validator 逐条验收，失败则在预算内重试。全部子任务完成后，meso 层运行集成验证；若集成失败，系统对失败做语义分类并选择最小修复范围——例如判定为“接口不匹配”则仅局部重跑受影响的子任务，判定为“设计缺陷”则回退 macro 层重新拆解。整个过程的进度持续落盘，用户可随时在终端旁观代理的实时输出，也可在每个子任务后暂停审查。

\subsection{与人工干预点的关系}

上述两个案例共同体现了 WhatyTerm 的一个设计取向：\emph{自动化不追求消灭人工干预，而是将人工干预点显式化、最小化并可配置}。哪些阶段可全自动、哪些阶段必须人工确认，由领域策略的阶段配置决定；这既保证了长程任务的自主推进，又在关键节点保留了必要的人类监督。
